\section{Analysis}
\label{sec:analysis}

\subsection{Predictor response and marginal utility}

The router observes how a candidate changes the frozen predictor's output.
The following results quantify how much that response tells us about the
candidate's marginal utility. Proofs are in \cref{app:proofs}.

\begin{proposition}[Response expansion]
\label{prop:expansion}
Assume that, for every episode, $f\mapsto\loss(f,y)$ is twice continuously
differentiable on the segment between $f_A(x)$ and $\fj(x)$, and
that $\|\nabla^2\loss(z,y)\|_2\leq L$ on that segment. With
$\dAj(x)=\fj(x)-f_A(x)$,
\begin{align}
 \marginal(j\mid A)
 &=
 -\E\!\left[\nabla\loss(f_A,y)^\top \dAj\right]
 -\tfrac12\E\!\left[\dAj^\top\nabla^2\loss(\xi,y)\dAj\right]
 -\lambda c_j,
 \label{eq:prop1-exact}\\
 \left|\marginal(j\mid A)
 +\E\!\left[\nabla\loss(f_A,y)^\top \dAj\right]+\lambda c_j\right|
 &\leq
 \frac{L}{2}\E\|\dAj\|^2 .
 \label{eq:prop1-bound}
\end{align}
\end{proposition}

The leading term is the alignment between the loss gradient at the current
prediction and the direction in which the candidate moves the predictor. The
remainder is second order in the response magnitude, so the first-order term is
accurate exactly when the candidate perturbs the predictor little, and a large
response can cost at most $(L/2)\E\|\dAj\|^2$.

\begin{proposition}[Exact identity for Bregman losses]
\label{prop:bregman}
Let $\loss(f,y)=D_\phi(y,f)=\phi(y)-\phi(f)-\nabla\phi(f)^\top(y-f)$ for convex
differentiable $\phi$. Then
\begin{equation}
 \marginal(j\mid A)
 =
 \E\!\left[\left(\nabla\phi(\fj)-\nabla\phi(f_A)\right)^\top(y-f_A)\right]
 -\E D_\phi\!\left(f_A,\fj\right)
 -\lambda c_j .
 \label{eq:prop2}
\end{equation}
\end{proposition}

Equation~\eqref{eq:prop2} is exact and requires no smoothness constant. The
second term is a non-positive movement penalty in the Bregman geometry. With
$\phi(z)=\|z\|^2$ it reduces to the squared-loss identity of
\cref{eq:squared-marginal}. If $\phi$ is $\mu$-strongly convex and $L$-smooth,
the penalty is bounded by $(\mu/2)\|d\|^2$ and $(L/2)\|d\|^2$ from below and
above.

\begin{proposition}[Cross-entropy, explicit constants]
\label{prop:ce}
For logits $f\in\R^C$, one-hot $y$, and softmax cross-entropy
$\loss(f,y)=-f_y+\log\sum_k e^{f_k}$ with $p=\softmax(f)$, the Hessian
$\nabla^2\loss=\diag(p)-pp^\top$ is positive semidefinite with spectral norm at
most $1/2$. Consequently
\begin{equation}
 \langle e_y-p_A,\dAj\rangle-\tfrac14\E\|\dAj\|^2-\lambda c_j
 \;\leq\;
 \marginal(j\mid A)
 \;\leq\;
 \langle e_y-p_A,\dAj\rangle-\lambda c_j .
 \label{eq:prop3}
\end{equation}
\end{proposition}

The upper bound holds because the second-order term penalises any positive
semidefinite curvature; the lower bound says that a large response can lose at
most a quarter of its squared size in logit space.

\begin{proposition}[Response sufficiency for squared loss]
\label{prop:suff}
For squared loss, conditioning on the state gives
\begin{equation}
 \E\!\left[\widetilde{\marginal}(j\mid A)\mid x,A\right]
 =
 \langle g_A(x),\dAj(x)\rangle-\|\dAj(x)\|^2-\lambda c_j,
 \quad
 g_A(x)=2(\mu(x)-f_A(x)),
 \label{eq:prop4}
\end{equation}
where $\mu(x)=\E[y\mid x,A]$.
\end{proposition}

Thus, conditional on the state, the Bayes-optimal marginal-utility predictor is
affine in the response vector plus a known quadratic penalty. This is the
parameterisation used by the response-aware router: state features supply the
coefficient $g_A(x)$, response features supply $\dAj$, and the quadratic
term is a function of the response alone. Any score that is independent of the
response can only recover the state average and provably omits the
$\langle g_A,\dAj\rangle$ interaction that the headroom diagnostic of
\cref{eq:headroom} measures. For a general differentiable loss the same
argument replaces $2(f_A-y)$ by $\nabla\loss(f_A,y)$, which is not a function
of the residual alone, so both state and response must be observed; this is
why the router head receives both. The bounds are conditional on the state and
require the loss gradient at the current prediction; at routing time the label
is unobserved and the router estimates the expectation of the first-order term
from training episodes. The theory therefore explains the feature design and
the second-order error of a response-based surrogate rather than providing a
label-free guarantee.

\subsection{From estimation error to routing regret}

\begin{theorem}[One-step routing regret]
\label{thm:onestep}
Suppose that, at a fixed state $A$, the utility estimator satisfies
$|\widehat{\marginal}(j\mid A)-\marginal(j\mid A)|\leq\epsilon$ for every
remaining candidate. If $\widehat{j}$ maximizes the estimated utility and
$j^*$ maximizes the true utility, then
\begin{equation}
 \marginal(j^*\mid A)-\marginal(\widehat{j}\mid A)\leq 2\epsilon.
\end{equation}
\end{theorem}

\begin{proof}
By the error bound,
$\marginal(j^*\mid A)\leq\widehat{\marginal}(j^*\mid A)+\epsilon$.
The maximization rule gives
$\widehat{\marginal}(j^*\mid A)\leq
\widehat{\marginal}(\widehat{j}\mid A)$, and the error bound applies once more.
\end{proof}

Combining \cref{thm:onestep} with \cref{prop:expansion} bounds the regret of a
response-based router by twice the estimation error of its first-order
surrogate plus the second-order response term. The theorem controls one
routing step. A sequence of locally accurate choices can still visit a
different state from exact greedy selection. When the induced set utility is
normalized, monotone, and submodular, the standard approximate-greedy argument
extends the local bound to
\begin{equation}
 F(A_{\mathrm{R-MUR}})
 \geq
 \left(1-\left(1-\frac{1}{B}\right)^B\right)F(A^*)-2B\epsilon,
\end{equation}
where $B$ is the context budget. The one-step result itself does not require
monotonicity or submodularity.
